% =====================================================================
% DRAFT REPLACEMENT for manuscript.tex Section 3 ("The doubletree algorithm")
% Reframed onto the working-model ATT theta*_n + impossibility spine.
%
% STATUS: FOR REVIEW. Not yet merged into manuscript.tex.
% Source of truth: notes/single-tree-inference.tex (theory) +
%   quality_reports/specs/2026-08-04_estimand-decision-theta-star.md (decision).
% Notation: uses manuscript macros (common-defs.tex). Drops in as-is.
%
% This block is written to REPLACE current manuscript.tex lines 135-165
% (\section{The doubletree algorithm} ... end of the epsilon_n caveat remark).
% It keeps \label{sec:rashomon} so existing \ref{sec:rashomon} cross-refs survive.
% The old Rashomon-intersection material moves to the appendix (Appendix B already
% exists); §3 now presents Rashomon as ONE way to satisfy the recovery event.
%
% Dependencies assumed present in the manuscript when merged:
%   - Lemma 1 (working-model ATT bilinear/DR bias) -- import from
%     notes/lemma1-bias-proof.tex, referenced here as \ref{lem:wm-bias}.
%   - Prop 1 (tree nuisance rate, §2) -- referenced as \ref{prop:tree-rate}.
%   - Cross-fit estimator thetahat_crossfit (§2) -- the delta-hat anchor.
%   - CSA citation key: \citet{chen2022dml} (ADD to .bib -- see note at end).
% =====================================================================

\section{Single interpretable tree per nuisance}
\label{sec:rashomon}

The cross-fit estimator of Section~\ref{sec:methods} is valid but produces $K$
distinct trees per nuisance---one per fold---which is precisely the interpretability
cost we set out to remove: a stakeholder auditing the propensity model must reconcile
$K$ different decision rules rather than inspect one. Our goal in this section is a
\emph{single} interpretable tree per nuisance, with leaf values fit on \emph{all} $n$
observations, together with valid $\sqrt{n}$ inference and \emph{no sample-splitting
at the inference stage}.

Achieving this forces a change in estimand, and that change is the intellectual core
of the paper. We show (Theorem~\ref{thm:wm-clt}) that a fixed, bounded-complexity tree
admits valid full-sample inference for the \emph{working-model ATT} $\thetastar$---the
ATT of the displayed tree model---via leave-one-out (LOO) algorithmic stability
\citep{chen2022dml} rather than cross-fitting. We then show
(Proposition~\ref{prop:gap}, Corollary~\ref{cor:impossible}) that this same
interpretability constraint makes the \emph{true} ATT $\theta_0$ generically
\emph{unreachable} at the $\sqrt{n}$ rate: no fixed-complexity interpretable tree can
deliver a $\sqrt{n}$-valid confidence interval for $\theta_0$ under H\"older smoothness.
The gap $\thetastar - \theta_0$ is reported as an honest diagnostic $\deltahat$, never
used to correct the estimate. A structural piecewise-constant assumption recovers
$\theta_0$ as a special case (Corollary~\ref{cor:structural}). This is a sharp
statement of the price of interpretability, and it is the paper's spine.

\subsection{The working-model ATT}
\label{sec:wm-estimand}

For any candidate nuisances $(e, f) \in \Ltwo(\Prob)$ plugged into the score
\eqref{eq:score}, the \emph{working-model ATT} $\theta(e, f)$ is the solution of the
population moment condition $\E_\Prob[\psi(\bO; \theta, e, f)] = 0$; at the truth,
$\theta(e_0, \mu_0) = \theta_0$. The following bias identity is the workhorse
(proved as Lemma~\ref{lem:wm-bias}; the bilinear, doubly robust structure is what
makes the tree-approximation error enter as a \emph{product}).

\begin{lemma}[Working-model ATT bias]
\label{lem:wm-bias}
Under Assumption~\ref{assump:causal-inf-assumps} (strong positivity with constant
$c$), for any $(e, f)$ with $e \in [c, 1-c]$ a.s.,
\begin{equation}
\label{eq:wm-bias}
\theta(e, f) - \theta_0
  = \frac{1}{\pi}\, \E_\Prob\!\left[ \frac{e_0 - e}{1 - e}\,(\mu_0 - f) \right],
\qquad
\abs{\theta(e, f) - \theta_0}
  \le \frac{1}{\pi c}\, \twonorm{e - e_0}\, \twonorm{f - \mu_0}.
\end{equation}
\end{lemma}

A partition $\tau$ of $\cX$ into $L = \abs{\tau}$ axis-aligned leaves induces
piecewise-constant nuisances by cell-wise population means. On the adaptive
discretization grid $G_n$ of Section~\ref{sec:discretization} (quantile bins
$b_n \asymp \log n$ per feature), let $\tau_0 = \tau_0(G_n)$ denote the
\emph{grid-optimal} partition of bounded complexity $L = o(\sqrt{n})$, with
piecewise-constant nuisances $e_{\tau_0}, \mu_{\tau_0}$. We define the
\textbf{interpretable working-model ATT}
\begin{equation}
\label{eq:thetastar}
\thetastar \;:=\; \theta(e_{\tau_0}, \mu_{\tau_0}),
\end{equation}
the ATT functional evaluated at the grid-optimal tree model. This---not $\theta_0$---is
the target of the inference below.

\begin{remark}[Estimand legitimacy]
\label{rem:estimand-legit}
$\thetastar$ is a projection / working-model parameter in the established tradition of
assumption-lean and structure-agnostic regression targets. It is the ATT a practitioner
would report if they believed the displayed tree correctly described the nuisance
functions. What is new here is that (i) the working model is a \emph{data-displayed
interpretable tree}---the estimator is the deliverable---and (ii) inference is valid on
all $n$ with no cross-fitting, via LOO stability. We are careful never to call
$\thetastar$ ``the ATT'' without the working-model qualifier, and we never claim the
interval covers $\theta_0$; the relationship between the two is characterized exactly
in Section~\ref{sec:gap}.
\end{remark}

\subsection{The estimator}
\label{sec:wm-estimator}

Given a data-selected partition $\tauhat$ (from the cross-fit trees of
Section~\ref{sec:methods} or a Rashomon/modal menu; see Section~\ref{sec:selection}),
we \emph{fix} its structure and refit leaf values as within-leaf sample means on all
$n$ observations:
\[
\ehat_{\tauhat}(\bx) = \Prob_n[A \mid \bX \in \ell(\bx)],
\qquad
\muhat_{\tauhat}(\bx) = \Prob_n[Y \mid \bX \in \ell(\bx),\, A = 0],
\]
where $\ell(\bx)$ is the leaf containing $\bx$ and $\Prob_n$ denotes the empirical
mean. The full-sample plug-in $\thetahat$ solves
$\Prob_n[\psi(\bO; \theta, \ehat_{\tauhat}, \muhat_{\tauhat})] = 0$. The displayed
deliverable is exactly these two trees plus $\thetahat$ and its interval.

We impose the following complexity-control condition on the selected partition, in
addition to the causal assumptions of Section~\ref{sec:methods}.

\begin{assumption}[Balanced bounded complexity]
\label{assump:leaves}
The selected partition $\tauhat$ has $L = \abs{\tauhat}$ leaves with
$L = o(\sqrt{n})$, and each leaf count satisfies $n_\ell \ge \underline{c}\, n / L$
for a constant $\underline{c} > 0$ (a minimum-leaf-size regularization). Consequently
$\min_\ell n_\ell \gtrsim n / L \gg \sqrt{n}$.
\end{assumption}

\subsection{Full-sample inference for the working-model ATT}
\label{sec:wm-clt}

The score \eqref{eq:score} is linear in $\theta$: writing
$m(\bO; \theta, g) = a(\bO; g)\,\theta + \nu(\bO; g)$ with nuisance $g = (e, f)$, we
have $a(\bO; g) = -A/\pi$ (free of $g$) and all nuisance dependence confined to
$\nu$. This places the estimator in the linear-moment framework of
\citet{chen2022dml}, whose central insight is that an estimator built from an
\emph{algorithmically stable} nuisance map is asymptotically equivalent to its
cross-fit counterpart---so sample-splitting can be dropped at the inference stage.
The next lemma establishes the required stability for a \emph{fixed} partition; its
proof is in Appendix~\ref{app:single-tree}.

\begin{lemma}[Leaf-refit LOO stability]
\label{lem:leaf-stability}
Fix a partition $\tauhat$. Under Assumptions~\ref{assump:causal-inf-assumps} and
\ref{assump:leaves}, replacing a single observation changes each leaf mean by
$O(L / n)$ uniformly, and the induced perturbation of $\nu$ satisfies the
\citet{chen2022dml} stability conditions at rate $O(L / n) = o(n^{-1/2})$, together
with a mean-squared-continuity condition. Hence the stochastic-equicontinuity
condition of \citet{chen2022dml} holds for the moment \eqref{eq:score}.
\end{lemma}

Structure-selection instability is quarantined by a \emph{partition-fixing event}.

\begin{assumption}[Partition-fixing event]
\label{assump:En}
Let $E_n = \{\tauhat = \tau_0\}$ be the event that the selector recovers the
grid-optimal partition, and require the failure probability to vanish faster than the
root-$n$ rate: $\Prob(E_n^c) = o(n^{-1/2})$. This is strictly stronger than
$\Prob(E_n) \to 1$ (see Remark~\ref{rem:rate-necessity}) and is discharged by either
mechanism of Section~\ref{sec:selection}.
\end{assumption}

\begin{theorem}[Working-model CLT]
\label{thm:wm-clt}
Suppose Assumptions~\ref{assump:causal-inf-assumps}, \ref{assump:regularity},
\ref{assump:leaves}, and \ref{assump:En} hold, and the tree nuisance rates satisfy
$\E_{\bX}\twonorm{\hat g(\bX) - g_0(\bX)}^2 = o_p(n^{-1/2})$ (equivalently each
nuisance $o_p(n^{-1/4})$, delivered by Proposition~\ref{prop:tree-rate} when
$\alphabar > s/2$). Then
\begin{equation}
\label{eq:wm-clt}
\sqrt{n}\,\bigl(\thetahat - \thetastar\bigr) \leadsto N(0, V),
\qquad
V = \E\bigl[\psi(\bO; \thetastar, e_{\tau_0}, \mu_{\tau_0})^2\bigr],
\end{equation}
and a LOO-stability plug-in $\hat V \to_p V$. Hence
$\thetahat \pm z_{1-\alpha/2}\sqrt{\hat V / n}$ is an asymptotically valid confidence
interval for $\thetastar$.
\end{theorem}

The proof (Appendix~\ref{app:single-tree}) verifies the \citet{chen2022dml}
conditions: Neyman orthogonality of \eqref{eq:score} and second-order smoothness are
supplied by Lemma~\ref{lem:wm-bias} (the bilinear bias is the second-order remainder);
stochastic equicontinuity is Lemma~\ref{lem:leaf-stability} on the fixed $\tau_0$; and
the recovery event passes the conditional CLT to the unconditional law because
$\sqrt{n}\,\Prob(E_n^c) = o(1)$ under Assumption~\ref{assump:En}.

\begin{remark}[Why $o(n^{-1/2})$, not $o(1)$, is necessary]
\label{rem:rate-necessity}
If $\Prob(E_n^c)$ vanished only slower than $n^{-1/2}$ (e.g.\ $\asymp n^{-1/4}$), the
recovery-failure term would contribute $\sqrt{n}\,\Prob(E_n^c) \to \infty$ to the
centering, breaking the CLT even though structure recovery ``succeeds with probability
tending to one.'' This is the post-selection subtlety: consistency of selection is
insufficient; the failure \emph{rate} must beat the estimator's own $\sqrt{n}$ rate.
It is also why selection cannot be certified by simulation alone---a $2$--$5\%$
empirical misrecovery rate looks benign, but its rate in $n$ is what governs validity,
and that rate is invisible at moderate $n$. We disclose this proactively as motivation
for the held-out selection route below.
\end{remark}

\subsection{The gap to \texorpdfstring{$\theta_0$}{theta0} and an impossibility}
\label{sec:gap}

The working-model CLT is valid for $\thetastar$. How far is $\thetastar$ from the true
ATT $\theta_0$? Under smoothness, the answer is: too far to close at the $\sqrt{n}$
rate.

\begin{proposition}[Discretization gap]
\label{prop:gap}
Under a H\"older-$\beta$ smoothness class for $(e_0, \mu_0)$ and the adaptive grid
$G_n$ ($b_n \asymp \log n$ bins per feature),
\[
\thetastar - \theta_0 = O\!\bigl((\log n)^{-2\beta}\bigr),
\qquad\text{hence}\qquad
\sqrt{n}\,\bigl(\thetastar - \theta_0\bigr) \to \infty
\ \text{ for every fixed } \beta > 0.
\]
\end{proposition}

\begin{corollary}[Impossibility for $\theta_0$]
\label{cor:impossible}
No estimator built from a single \emph{fixed}, \emph{bounded-complexity}
($L = o(\sqrt{n})$) interpretable tree per nuisance on the adaptive grid admits a
$\sqrt{n}$-valid confidence interval for $\theta_0$ under H\"older-$\beta$ smoothness.
Driving the gap to $o(n^{-1/2})$ requires polynomial grid refinement, which forces
$L \ne o(\sqrt{n})$ (breaking Lemma~\ref{lem:leaf-stability}) or an off-tree debiasing
term (breaking ``the tree is the estimator'').
\end{corollary}

Corollary~\ref{cor:impossible} is the paper's central negative result: among
\{fixed tree\}, \{interpretable\}, \{$\sqrt{n}$-honest for $\theta_0$\}, one may have
any two. It is sharper and more informative than any positive $\theta_0$ construction
built on off-tree corrections, and it is why we target $\thetastar$.

\begin{corollary}[Structural piecewise-constant special case]
\label{cor:structural}
If $(e_0, \mu_0)$ are \emph{exactly} piecewise-constant on some finite (unknown)
partition $\tau^\ast$---a structural, not smoothness, assumption---then for $n$ large
enough that $G_n$ resolves $\tau^\ast$ we have $e_{\tau_0} = e_0$,
$\mu_{\tau_0} = \mu_0$, hence $\thetastar = \theta_0$, and Theorem~\ref{thm:wm-clt}
delivers a $\sqrt{n}$-valid confidence interval for the \emph{true} ATT $\theta_0$.
\end{corollary}

This is the only coherent route to honest $\theta_0$ inference within the fixed-tree
constraint, and it is exactly the regime our binary-covariate implementation targets
(cf.\ the binary-rate remark of Section~\ref{sec:methods}), where a tree with enough
leaves represents any function on the active subcube exactly.

\begin{remark}[The diagnostic $\deltahat$ is not a correction]
\label{rem:delta}
Define $\deltahat = \thetahat - \thetahat_{\text{crossfit}}$, the difference between
the displayed-tree estimate and the cross-fit anchor of Section~\ref{sec:methods}. By
Lemma~\ref{lem:wm-bias}, $\deltahat$ is a plug-in estimate of the gap
$\thetastar - \theta_0$ (up to $O_p(n^{-1/2})$ noise) and is reported as an
\emph{honesty flag}. It is never subtracted from $\thetahat$: by
Proposition~\ref{prop:gap} the gap is asymptotically non-negligible relative to the CI
width, so a $\deltahat$-correction capable of recovering $\theta_0$ would already
constitute a solution of the $\theta_0$ problem, contradicting
Corollary~\ref{cor:impossible}. In practice we report $\deltahat$ and
$\deltahat / \widehat{\mathrm{SE}}$ alongside every estimate: small values indicate the
displayed tree is a faithful summary of the ATT; large values flag a meaningful gap
between the working model and the truth.
\end{remark}

\subsection{Discharging the partition-fixing event: structure selection}
\label{sec:selection}

It remains to supply a selector satisfying Assumption~\ref{assump:En}. Two routes are
available; the choice is orthogonal to the estimand (both target the same
$\thetastar$) and is settled by the empirical recovery behavior of the selector.

\paragraph{Route B (exact recovery).}
If the grid-optimal $\tau_0$ has a risk margin $\Delta_n > 0$ over its competitors, a
modal selector over $M_n$ splits obeys a union-bound tail
$\Prob(E_n^c) \le S_n \exp(-M_n \Delta_n)$, and Assumption~\ref{assump:En} holds
whenever $M_n \Delta_n - \log S_n \gg \tfrac12 \log n$. Under the log-$n$ grid this
requires $M_n \gg (\log n)^{1 + \kappa}$. This route keeps the pure ``single tree, all
$n$, no split anywhere'' story, but the margin $\Delta_n > 0$ is generic only in the
structural piecewise-constant regime of Corollary~\ref{cor:structural}; when the signal
is merely H\"older, the risk landscape is flat near the optimum and exact recovery can
fail. The \emph{Rashomon-intersection} construction (Appendix~\ref{app:rashomon}) is
one concrete mechanism that supplies such a margin, by selecting the structure common
to near-optimal trees across folds.

\paragraph{Route A (held-out selection).}
Partition the sample into a selection fold $\cS_n$ of size $\lfloor \rho_n n \rfloor$
and an inference fold $\cI_n$, with selection fraction $\rho_n \to 0$ and
$\rho_n n \to \infty$. Select $\tauhat$ using $\cS_n$ only, then refit leaf values on
$\cI_n$. Conditional on $\cS_n$ the partition is fixed, so
$\Prob(E_n^c \mid \cS_n) = 0$ trivially and Lemma~\ref{lem:leaf-stability} applies to
the fixed $\tauhat$ on $\cI_n$; since $\rho_n \to 0$, the asymptotic variance is
unchanged and Theorem~\ref{thm:wm-clt} holds at the full $\sqrt{n}$ rate
(Appendix~\ref{app:single-tree}). This route is robust to selection instability by
construction and requires no margin.

\begin{remark}[Honest scope of ``no sample-splitting'']
\label{rem:route-a-scope}
Route A reintroduces a split at the \emph{selection} stage---$\tauhat$ is chosen on
$\cS_n$---but preserves the property CSA buys: \emph{no split at the inference stage}.
Leaf values and the CLT use nearly the full sample $\cI_n$, with LOO stability standing
in for cross-fitting. The displayed tree is still a single interpretable tree; only its
\emph{structure} was chosen on a subsample, while its \emph{leaf values}---the numbers a
reader acts on---use nearly all the data. This is strictly weaker than the Route B ideal
but is the price of validity when the margin is flat and exact recovery fails.
\end{remark}

Which route to adopt is a direct read on the empirical structure-recovery frequency
$\widehat{\Prob}(E_n)$, reported in Section~\ref{sec:results}: Route B where recovery
$\to 1$ (the structural-PC regimes), Route A elsewhere. In either case the estimand and
Theorem~\ref{thm:wm-clt} are unchanged.


% =====================================================================
% MERGE / TODO NOTES (delete before final):
%
% 1. BIB: add CSA key to the .bib file:
%      @inproceedings{chen2022dml,
%        author = {Chen, Qizhao and Syrgkanis, Vasilis and Austern, Morgane},
%        title  = {Debiased Machine Learning without Sample-Splitting for
%                  Stable Estimators},
%        booktitle = {Advances in Neural Information Processing Systems (NeurIPS)},
%        year   = {2022}, note = {arXiv:2206.01825}}
%
% 2. APPENDIX: add \label{app:single-tree} (new appendix folding in
%    notes/single-tree-inference.tex: Lemma leaf-stability proof, Thm 1 proof,
%    Route A lemma proof). Re-label the existing Rashomon appendix
%    \label{app:rashomon} and demote it to "one mechanism for Route B."
%
% 3. LEMMA 1: import notes/lemma1-bias-proof.tex as \ref{lem:wm-bias} (currently
%    referenced but proof lives in the note). Reconcile the score convention
%    (manuscript eq:score already matches the note's empirical-treated-mean form).
%
% 4. thetahat_crossfit: ensure §2 names the cross-fit estimator
%    \thetahat_{\text{crossfit}} so \ref in rem:delta resolves (manuscript §2
%    currently calls it \thetahat_{\text{crossfit}} at line 125 -- OK).
%
% 5. CROSS-REFS INTO §4/§6: the outline's §4 (simulation) speaks of "coverage of
%    theta*" and "delta-hat/SE calibration" -- those align with this reframe.
%    §6 limitations should cite Corollary impossible as the headline limitation.
%
% 6. REMOVED from old §3: the epsilon_n = o(n^{-1/2}) catch-22 discussion and the
%    intersection-overhead corollary move wholesale to Appendix B (app:rashomon);
%    they survive as Route B mechanics, not main-text theory.
% =====================================================================
